% !TEX program = lualatex
% Compile with LuaLaTeX, like the thesis. The fonts in ./fonts/ are copies of those in the thesis:
% LuaLaTeX caches fonts by name with their relative path, so both documents must use ./fonts/.
\documentclass[11pt,a4paper]{article}

\usepackage[british]{babel}
\usepackage{fontspec}
\setmainfont{UGentPannoText.otf}[
    Path = ./fonts/,
    UprightFont = *-Normal,
    BoldFont    = *-SemiBold,
    ItalicFont = *-Normal,
    ItalicFeatures = {FakeSlant=0.2},
    BoldItalicFont = *-SemiBold,
    BoldItalicFeatures = {FakeSlant=0.2},
]
\usepackage{microtype}
\usepackage[margin=2.5cm]{geometry}
\usepackage{array,booktabs}
\usepackage[hidelinks,pdftitle={Summary of Changes},pdfauthor={Maxime Deforche}]{hyperref}

\setlength{\parindent}{0pt}
\setlength{\parskip}{0.5em}

% Heading, kept on the same page as the table that follows.
\newcommand{\heading}[1]{%
  \par\addvspace{1.5em}%
  {\raggedright\large\bfseries #1\par}\nopagebreak[4]%
  \vspace{0.2em}\nopagebreak[4]%
}

% One table per part of the dissertation: location on the left, description on the right.
\newlength{\loccol}
\setlength{\loccol}{3.7cm}
\newenvironment{changes}{%
  \begin{tabular}{@{}>{\raggedright\arraybackslash}p{\loccol}
                     >{\raggedright\arraybackslash}p{\dimexpr\textwidth-\loccol-2\tabcolsep\relax}@{}}
    \toprule
}{%
    \bottomrule
  \end{tabular}\par
}
\newcommand{\change}[2]{#1 & #2 \\ \addlinespace[0.45em]}
% Location: section or label, followed by the paragraph or item.
\newcommand{\loc}[2]{#1\enspace #2}
% Title of an unnumbered subsection, below the location.
\newcommand{\sub}[1]{\newline{\footnotesize\itshape #1}}
% Short label that opens each description.
\newcommand{\topic}[1]{\textbf{#1.}}

\begin{document}

\begin{center}
  {\LARGE\bfseries Summary of Changes\par}
  \vspace{0.4em}
  {\large Management and Analysis of Short Texts: A Graph-Based Approach\par}
  \vspace{0.8em}
  Maxime Deforche\par
  \today\par
\end{center}
\vspace{0.6em}

This document summarises the changes made to the dissertation after the internal defence, in response to the remarks of the jury members, together with some other larger improvements for text and print quality.
Locations refer to the revised version and provide a reference to the specific section and paragraph, e.g., 1.2.2\enspace ¶\,4 is the fourth paragraph of Section 1.2.2.
Unnumbered subsections are indicated with their title.

%======================================================================
\heading{General}

\begin{changes}
  \change{\loc{Line numbers}{}}{Removed the proof line numbers from the dissertation.}
  \change{\loc{Language}{}}{Corrected spelling and writing errors throughout the dissertation and improved the language in some sentences, following the suggestions of the jury members. (too many to list in detail)}
  \change{\loc{Terminology}{}}{Standardised some terms and wordings throughout the dissertation, to improve consistency and readability. (too many to list in detail)}
  \change{\loc{Images}{}}{Replaced some images with higher-resolution versions for better readability and print quality.}
  \change{\loc{Page size}{}}{Changed the dimensions of the thesis to 160×240\,mm, according to the guidelines of the faculty, and adjusted the margins accordingly.}
  \change{\loc{Layout}{}}{Some additional layout changes, such as moving some objects, change the size of some objects, ...}
  \change{\loc{Bibliography}{}}{Checked the urls in the bibliography to make sure that they are still valid and updated the ``visited on'' dates accordingly.}
  \change{\loc{Bibliography}{}}{Updated the formatting of the bibliography to ensure consistency and proper citation style.}
\end{changes}

%======================================================================
\heading{Preliminary Pages}

\begin{changes}
  \change{\loc{Dankwoord}{}}{\topic{Acknowledgements} Wrote the acknowledgements.}
  \change{\loc{Summary}{}}{\topic{Accessibility} Adapted the summary to be more accessible and readable, while still maintaining its technical nature. The content is equal to the original summary, but the language has been improved, some sentences have been reworded, and some clarifications have been added. (too many to list in detail)}
  \change{\loc{Samenvatting}{}}{\topic{Dutch summary} Made the same adaptations to the Dutch summary to match the updated English version. Also changed some translations slightly for a more accurate and natural Dutch translation. (too many to list in detail)}
\end{changes}

%======================================================================
\heading{Chapter 1: Introduction}

\begin{changes}
  \change{\loc{1.2.2}{¶\,4}}{\topic{Interpretability} Softened the claim that more complex models are generally less interpretable, by acknowledging that this is not always the case and that the claim in the thesis mostly refers to deep neural networks as black boxes.}
  \change{\loc{1.2.3}{¶\,5}}{\topic{Foundation models} Mentioned foundation models and that they can (partially) generalise over languages, but can have issues with low-resource languages, depending on the transferable features of the language and the availability of training data.}
  \change{\loc{1.3.1}{¶\,2--3}}{\topic{Epigram types} Made small additions to the introduction of epigram types, to emphasise that these are a single defensible scholarly interpretation of the corpus rather than an objective ground truth, similar to what is mentioned in the later chapters.}
  \change{\loc{1.3.2}{¶\,2}\ and\newline\loc{1.4.2}{¶\,2}}{\topic{Embeddings and NLP tools} Slightly softened the claim that embedding approaches and other NLP tools are not suitable for Byzantine texts at all. Acknowledged that some exist, but clarified that these are mostly based on normalised or edited texts, which is not the case for the DBBE corpus. Also mentioned that some embedding approaches aslo result from another subproject of the Interconnected Texts project.}
  \change{\loc{1.3.3}{¶\,1}}{\topic{Intertextuality} Explicitly mentioned the problem of intertextuality, what it is, and that this work can contribute to it.}
  \change{\loc{1.6}{item 3}}{\topic{Publication details} Added publication information (DOI, volume, \dots) for \textit{Verse, Variation, and Meaning: Annotating Semantic Distance in Greek Marginal Texts}, now that it has been published.}
  \change{\loc{1.6}{item 7}}{\topic{Publication details} Added publication information (DOI, volume, \dots) for \textit{Ranking L-grades for Veracity Handling in Logic Reasoning}, now that it has been published.}
\end{changes}

%======================================================================
\heading{Chapter 2: A Graph Data Model for Short Interconnected Texts}

\begin{changes}
  \change{\loc{Opening}{}}{\topic{Epigraph} Changed the epigraph to a quote by George E.\ P.\ Box, instead of the quote by Abby Hoffman, which did not really fit the thesis.}
  \change{\loc{2.3.1}{¶\,1}}{\topic{Indexing} Fixed the different indexing schemes in this section: the characters of a text and its substrings are now both indexed from 0 instead of one being indexed from 1 and the other from 0.}
  \change{\loc{2.6.1}{¶\,2}\ and\newline\loc{2.7}{¶\,4}}{\topic{Spelling shifts} Added a short part about other shifts in the Greek language beyond itacism, to clarify the impact of dealing with them on the graph size and to consider them in future work.}
\end{changes}

%======================================================================
\heading{Chapter 3: A Hierarchical Orthographic Similarity Measure for Short Texts}

\begin{changes}
  \change{\loc{3.5.1}{¶\,2}}{\topic{Expert groupings} Clarified the limitations of the expert-curated groupings, similar to the discussions in later chapters: they represent one defensible scholarly interpretation rather than a perfect ground truth.}
  \change{\loc{3.5.2}{¶\,3}\sub{Two-Level String Matcher}}{\topic{Parameter $\alpha_Q$} Clarified that the parameter $\alpha_Q$ (In earlier version of the thesis, also noted as just $\alpha$) of the Two-Level string matcher differs from the threshold $\alpha$ of the $\alpha$-cut in the validation process.}
  \change{\loc{3.5.4}{Figure 3.9}\sub{Query-Based Grouping of Similar Text Nodes}}{\topic{Figure labels} Fixed the epigram id labels in Figure 3.9, which were incorrectly swapped between epigrams 22826 and 21182. Also corrected the follow-up discussion accordingly, including the groups and word positions of these two epigrams. This change should resolve the confusion about the results of this image, as mentioned in one of the readering reports.}
\end{changes}

%======================================================================
\heading{Chapter 4: A Context-Aware and Responsible Similarity Measure for Short Texts}

\begin{changes}
  \change{\loc{Opening}{}}{\topic{Epigraph} Changed the attribution of the epigraph to mention that it is a popular paraphrase of a quote from Aristotle, rather than a direct quote.}
  \change{\loc{4.4.2}{Figure 4.1}}{\topic{Figure caption} Extended the caption to explain the aggregated similarity $s$ and the confidence score $c$ shown in the figure.}
  \change{\loc{4.5.4}{¶\,2, 5}\sub{Quality Measure Results}}{\topic{Precision and recall} Added some more fine-grained discussion about the different trends in precision and recall in the experiments of this chapter, such as how the metadata raise recall across most thresholds. This way not all focuss is put on the $F_1$-score.}
\end{changes}

%======================================================================
\heading{Chapter 5: Ranking of Satisfaction-Confidence Pairs}

\begin{changes}
  \change{\loc{Opening}{}}{\topic{Publication details} Added publication information (DOI, volume, \dots) for \textit{Ranking L-grades for Veracity Handling in Logic Reasoning}, now that it has been published.}
  \change{\loc{5.3.2}{¶\,5}}{\topic{Injective} Briefly explained ``injective'' where it is mentioned in the context of preorders, as there was no clear definition provided.}
  \change{\loc{5.4.1}{¶\,1}\ and\newline\loc{5.4.2}{¶\,2}}{\topic{Tie relation} Introduced a separate tie relation $\sim$ for equally preferred L-grades, instead of using equality with sub- and super-scripts.}
\end{changes}

%======================================================================
\heading{Chapter 6: Context-Aware and Responsible Short Text Clustering}

\begin{changes}
  \change{\loc{6.5.2}{¶\,3}\sub{Similarity Inputs and Weight Vector Configurations}}{\topic{Weight vector notation} Clarified the notation of the optimised weight vectors and added an example to prevent any confusion.}
  \change{\loc{6.5.2}{¶\,3--4}\sub{Quality Measures, Adjusted Mutual Information (AMI)}}{\topic{ARI and AMI} Made some additions to clarify the definitions and differences between the Adjusted Rand Index (ARI) and the Adjusted Mutual Information (AMI). Also added a small example to illustrate these differences and explained how the two measures relate to each other when one is higher or lower than the other.}
  \change{\loc{6.5.3}{¶\,3}\sub{Context-Aware Similarity Clustering}}{\topic{Precision and recall} Added some more fine-grained discussion about the different trends in precision and recall instead of mainly focusing on the $F_1$-score: the metadata raise recall and lower precision for all three algorithms.}
\end{changes}

%======================================================================
\heading{Chapter 7: Conclusions and Future Work}

\begin{changes}
  \change{\loc{7.2.2}{¶\,2}}{\topic{Embeddings and NLP tools} Slightly softened the claim that embedding approaches and similar techniques are not suitable for Byzantine texts at all, and mention the work done by the NLP subproject in the Interconnected Texts project. Also provide some future work to incorporate these techniques in the proposed methods.}
  \change{\loc{7.2.3}{¶\,5}}{\topic{Precision and recall} Reflected the more fine-grained discussion about precision and recall of chapters 4 and 6 in the conclusions, instead of mainly focusing on the $F_1$-score.}
  \change{\loc{7.5}{¶\,2}}{\topic{Spelling shifts} Added small remarks about the potential inclusion of other spelling shifts beyond itacism in future work, and the potential advantages of such an approach.}
  \change{\loc{7.5}{¶\,3}}{\topic{Outlier detection} Added some ideas about outlier detection as future work, for instance to identify epigrams that fail to connect to the rest of their expert-curated group.}
  \change{\loc{7.5}{¶\,3}}{\topic{PageRank} Mentioned more potential analysis techniques, including but not limited to PageRank, in future work.}
  \change{\loc{7.5}{¶\,4}}{\topic{Generalisation and complexity} Made it even more explicit that experiments on other corpora are needed to validate the generalisability of the proposed methods, and described the effects on the computational complexity of the current approach on much larger corpora.}
  \change{\loc{7.5}{¶\,4}}{\topic{Genre and register} Mentioned the need for an analysis of the genre and register of the texts in future work, to better understand the impact of these factors on the proposed methods.}
  \change{\loc{7.5}{¶\,4}}{\topic{Intertextuality} Mentioned the potential of these methods for the general problem of intertextuality.}
\end{changes}

\end{document}
